%!TEX root = ../graduate-work.tex
\section{Conclusion}

\subsection{Summary of Results}

This thesis addresses the adaptation of parallelism techniques from large
language model training to the problem of formal neural network verification.
Two approaches were implemented and experimentally validated: Tensor Parallelism
(TP) and Fully Sharded Data Parallelism (FSDP), with FSDP extended to complete
verification via $\beta$-CROWN and Branch-and-Bound.

The following results were obtained:

\begin{enumerate}
    \item \textbf{Near-linear memory reduction with TP.}
    TP reduces peak GPU memory consumption proportionally to the number of GPUs
    (factor $\approx\!1.985$ at $P = 2$) by sharding both weight matrices and
    CROWN $A$-matrices. Specialised operators
    \texttt{BoundLinearTP\_Col}/\texttt{Row} with correct \texttt{AllReduce}
    and computational-graph partitioning were implemented.
    Bounds satisfy soundness for all tested models, but multi-zone models
    exhibit degrading bound tightness due to forced IBP substitution for
    intermediate bounds.

    \item \textbf{Bitwise-accurate bounds with FSDP.}
    FSDP shards weight matrices, assembling them layer-wise via
    \texttt{AllGather}. All tested configurations (MLPs of various widths and
    depths, ONNX models from VNN-COMP) produced bounds \textit{bitwise identical}
    to the single-GPU baseline ($\Delta lb = \Delta ub = 0.00$). Baseline
    memory (weight storage) is reduced by exactly 50\% (factor $1/P$ at $P{=}2$,
    since each rank stores exactly half of every weight tensor); peak memory
    during \texttt{compute\_bounds} by 17--39\% depending on the ratio of
    weights to $A$-matrices.

    \item \textbf{Automatic sharding.}
    Automatic sharding functions (\texttt{tp\_shard\_bounded\_module},
    \texttt{fsdp\_shard\_bounded\_module}) were implemented for both approaches;
    they transform any \texttt{BoundedModule} (including ONNX-loaded models)
    without manual network rewriting.

    \item \textbf{FSDP in complete verification ($\beta$-CROWN+BaB).}
    A fair comparison (identical batch size, round count, and the
    \texttt{force\_synchronous} flag) showed that FSDP adds
    $\leq\!5$~MB AllGather overhead and does not alter the complete-verification
    workload. On small models (${\sim}1.5$~MB weights) no memory savings are
    observed---the bottleneck is alpha tensors ($\gg\!1$~GB), not weights.
    FSDP preserves bound accuracy and scales to arbitrary BaB problems.

    \item \textbf{FSDP for convolutional layers (\texttt{BoundConv}).}
    Sharding was adapted to \texttt{BoundConv} (along the out\_channels axis).
    Validation on CIFAR-100 ResNet medium/large (VNN-COMP'24) yielded bitwise
    identical bounds. On ResNet-large (${\sim}95$~MB weights, $>\!4$~GB
    alpha tensors) the memory savings from weight sharding are outweighed by
    AllGather buffer overhead, confirming that alpha tensors dominate memory.

    \item \textbf{JIT compatibility and ViT verification.}
    Transformer verification required resolving JIT-compilation incompatibilities
    in \texttt{BoundReshape}, \texttt{BoundConcat}, and
    \texttt{BoundConstantOfShape}. After the fixes, the model \texttt{pgd\_2\_3\_16}
    from VNN-COMP'23 runs correctly in $\alpha,\beta$-CROWN with FSDP, producing
    bounds identical to the reference.
\end{enumerate}

In summary, TP and FSDP represent opposite ends of the \textit{memory--accuracy}
trade-off: TP achieves the highest memory savings (${\sim}2\times$) but loses
bound tightness as model depth grows; FSDP provides moderate peak memory savings
(17--39\% for incomplete verification) while preserving bitwise accuracy and
remaining compatible with $\beta$-CROWN+BaB.

\subsection{Future Work}

\textbf{Sharding alpha tensors.}
Experiments~6--8 show that in $\alpha$-CROWN+BaB mode the memory bottleneck is
per-neuron alpha tensors, not weight matrices. Distributing alpha tensors across
GPUs (analogous to Domain Parallelism) would eliminate this barrier and provide
meaningful memory savings during complete verification of large convolutional
networks.

\textbf{Alpha-CROWN optimisation with TP and FSDP.}
Gradient-based optimisation of alpha slopes is compatible with both approaches
since it uses the same backward pass as CROWN.

\textbf{Scaling to more GPUs.}
All experiments were conducted at $P = 2$. Validation at $P = 4, 8$ would
empirically characterise communication overhead and identify scalability limits.
FSDP at $P = 4$ theoretically achieves ${\sim}93\%$ baseline memory savings.

\textbf{Domain-Parallel BaB with input-split branching.}
The current FSDP--BaB integration operates in synchronous mode (all GPUs
process one sub-tree). Input-split branching---parallel exploration of different
sub-domains across GPUs---would enable not only memory savings but also genuine
verification speedup.
